%% Türkçe makale taslağı — trim mimarisi
%% Derleme: pdflatex makale_trim_tr.tex  (elsarticle.cls gerektirir)
%% Grafikler için yer tutucular [SEKIL X] biçiminde işaretlenmiştir.
%% Kaynak veriler: test_b_* serisi, test_orbit_* serisi,
%%   test_symm_vs_antisym, test_symm_basis_fit, test_radial_spin
%%   (ayrıntılar: false_edm_harmonic_sinir.md §12, trim_yontemi_pedagojik.md)

\documentclass[preprint,12pt]{elsarticle}

\usepackage[turkish]{babel}
\usepackage{amsmath,amssymb}
\usepackage{graphicx}
\usepackage{booktabs}
\usepackage{siunitx}
\usepackage{xcolor}
\newcommand{\todo}[1]{{\color{red}\textbf{(YAPILACAK: #1)}}}

\journal{Physical Review Accelerators and Beams}

\begin{document}

\begin{frontmatter}

\title{Dondurulmuş-spin proton depolama halkasında sahte EDM'nin
kaynağında bastırılması: EDM-kör yörünge trimi ve spin-sürülü
ölç--trimle döngüsünden oluşan iki kademeli mimari}

\author[inst1]{S.~Hac\i\"omero\u{g}lu\corref{cor1}}
\cortext[cor1]{Sorumlu yazar}
\ead{selcuk.haciomeroglu@gmail.com}

\address[inst1]{\todo{kurum bilgisi}}

\begin{abstract}
Proton elektrik dipol momenti (pEDM) dondurulmuş-spin deneyinde kuadrupol
mıknatısların düşey hizalama hataları, radyal manyetik alan üzerinden
gerçek EDM sinyalini taklit eden bir spin presesyonu (sahte EDM) üretir;
tolerans $\sim\SI{10}{\micro\meter}$ mertebesindedir. Yerleşik yaklaşım
hizalamayı \emph{ölçmek} (k-modülasyonu, ışın-tabanlı hizalama) iken, bu
çalışmada hizalama vektörünü hiç bilmeden sahte EDM'yi \emph{kaynağında
bastıran} iki kademeli bir trim mimarisi sunulmakta ve uçtan uca
simülasyonla doğrulanmaktadır. Birinci kademe, ışın konum monitörü (BPM)
okumalarından Fourier mod genliklerini kestirip mekanik kaydırıcılarla
geri besleyen, yapısal olarak EDM'ye kör bir yörünge trimidir; hangi
modların fit edilebileceği, ölçülen yörünge kazançlarının
$G_k=C/|Q_\mathrm{eff}^2-k^2|$ yasasına oturtulmasıyla kapalı biçimde
türetilen $k_\mathrm{max}^2<Q_\mathrm{eff}^2+C\,\sigma_q/\sigma_b$
eşiğiyle belirlenir; yasa, fit aralığının dışındaki $k=7..12$
modlarında \emph{öngörü} olarak $\le\%1$ doğrulukla sınanmıştır. İkinci
kademe, sahte EDM'nin kendisini gözlenebilir alan spin-sürülü
ölç--trimle döngüsüdür: spin--mod kuplaj katsayıları $c_k$'nin desenden,
fırlatma koşulundan ve genlikten bağımsızlığı gösterilmiş; rastgele
fazlı gerçekçi desenlerde çift-kuadratür kalibrasyonla
$\ge10^{7}\times$ bastırma elde edilmiştir. İki kademenin gerekliliği
dik-izdüşüm ayrıştırmasıyla kanıtlanmaktadır: 48 serbestlik derecesinin
23'ü (QF ve QD'nin aynı yönlü, ``simetrik'' kombinasyonları) zıt
gradyan işaretleri nedeniyle yörüngede $\sim2.6\times$, spinde
$\sim12\times$ bastırılmış fakat sıfır olmayan iz bırakır; bu içerik
BPM trimini $\sim10^{-4}$~rad/s tabanına demirlerken yalnız spin
kademesi tarafından görülür ve temizlenebilir. Mimari, CW/CCW ve
kuadrupol-çevirme iptal zincirinin giriş gereksinimi olan
$10^{-5}$~rad/s düzeyine giden bütünsel bir yol sunar.
\end{abstract}

\begin{keyword}
elektrik dipol momenti \sep depolama halkası \sep dondurulmuş spin \sep
hizalama hataları \sep kapalı yörünge bozulması \sep sistematik hatalar
\end{keyword}

\end{frontmatter}

%==============================================================
\section{Giriş}
\label{sec:giris}

Dondurulmuş-spin proton EDM deneyi, sihirli momentumdaki
($p=\SI{0.7007}{GeV/c}$) protonların tamamı-elektrik bir depolama
halkasında boylamsal polarizasyonla saklanmasına dayanır: manyetik
dipol momentin yatay düzlemdeki presesyonu durur ve gerçek bir EDM,
spini düşey eksene doğru saatler mertebesinde biriken bir hızla
döndürür~\cite{todo-pedm}. Bu hassasiyetin bedeli, aynı kanala
($dS_y/dt$) sızan her sistematik etkinin gerçek sinyalden ayırt
edilemez olmasıdır. En tehlikeli sistematik, kuadrupol mıknatısların
düşey hizalama hatalarından doğar: $\delta y$ kaçıklığı taşıyan bir
kuadrupol, kapalı yörüngeyi bozarak parçacığa net radyal manyetik alan
$B_x$ gösterir; Thomas--BMT denklemi üzerinden bu alan tam da EDM'nin
döndürdüğü eksende bir presesyon --- \emph{sahte EDM} --- üretir.
Omarov ve ark.~\cite{todo-omarov} bu mekanizma için hizalama
toleransını $\sim\SI{10}{\micro\meter}$ olarak vermiş; CW/CCW çift yönlü
saklama ve kuadrupol-çevirme (quad-flip) iptal teknikleriyle
$10^{-5}$~rad/s'lik bir artığın $10^{-9}$~rad/s altına
indirilebileceğini göstermiştir. Dolayısıyla hizalama altyapısının
deneye teslim etmesi gereken düzey $dS_y/dt\sim10^{-5}$~rad/s'dir.

Literatürdeki yaklaşımlar bu hedefe hizalamayı \emph{ölçerek} gider:
ışın-tabanlı hizalama (BBA), k-modülasyonu, LOCO türü tepki-matrisi
geri çatımları. Hepsinin ortak zorluğu, BPM elektronik ofsetlerinin
($\sigma_b\sim\SI{100}{\micro\meter}$) aranan sinyalle
($\sim\SI{10}{\micro\meter}$) arasındaki büyüklük uçurumudur. Bu
çalışmada farklı bir soru sorulmaktadır: \textbf{hizalama vektörünü hiç
bilmeden, sahte EDM'yi doğrudan gözleyip kaynağında --- kuadrupol
konumlarında --- bastırmak mümkün müdür?}

Yanıtın evet olduğu, iki tamamlayıcı geri-besleme kademesinin uçtan uca
simülasyonuyla gösterilmektedir:

\begin{enumerate}
\item \textbf{Yörünge kademesi (hızlı, EDM-kör):} BPM'lerle ölçülen
kapalı yörünge, kuadrupol kaçıklık modlarının kalibre edilmiş parmak
izlerine fit edilir; kestirilen genlikler mekanik kaydırıcılarla geri
beslenir. BPM spini görmediğinden bu kademe gerçek EDM sinyaline
yapısal olarak dokunamaz. Saniyeler içinde tamamlanır.
\item \textbf{Spin kademesi (yavaş, eksiksiz):} Sahte EDM'nin kendisi
($dS_y/dt$) gözlenebilir olarak kullanılır ve kalibre edilmiş
spin--mod kuplajları $c_k$ üzerinden trim hesaplanır. Bu kanal tüm
serbestlik derecelerini görür --- yörüngenin göremediği ``simetrik''
içerik dahil.
\end{enumerate}

Makalenin katkıları üç başlıkta toplanmaktadır. Birincisi, spin-sürülü
ölç--trimle döngüsünün çalışması için gereken üç önkoşulun ---
$c_k$ katsayılarının doğrusallığı, desenden bağımsızlığı ve fırlatma
koşulundan bağımsızlığı --- sayısal kanıtı
(Bölüm~\ref{sec:spin-trim}). İkincisi, yörünge kademesinin fit
aralığını belirleyen kapalı-form eşik teorisi ve bunun fit aralığı
dışında öngörü olarak doğrulanması (Bölüm~\ref{sec:orbit-trim}).
Üçüncüsü, iki kademenin neden birlikte gerektiğinin mekanizmalı
ispatı: kaçıklık uzayının dik ayrıştırmasıyla, BPM triminin tabanını
oluşturan ve yalnız spinin görebildiği simetrik alt-uzayın
nicelenmesi (Bölüm~\ref{sec:taban}). Bölüm~\ref{sec:elenen},
araştırma sürecinde denenip elenen alternatifleri mekanizmalarıyla
birlikte raporlar; Bölüm~\ref{sec:recete} deneysel reçeteyi verir.

%==============================================================
\section{Model ve simülasyon altyapısı}
\label{sec:model}

\subsection{Kafes ve izleyici}

Simüle edilen halka $R_0=\SI{95.49}{m}$ yarıçaplı, $n_\mathrm{FODO}=24$
hücreli (48 kuadrupol) bir tamamı-elektrik FODO kafesidir. Her hücre
ARC--drift--QF--drift--ARC--drift--QD--drift diziliminden oluşur;
deflektörler silindirik kapasitör ($n=1$ alan indisi), kuadrupoller
$L_q=\SI{0.4}{m}$ uzunluğunda ve $g_0=\SI{0.2}{T/m}$ gradyanındadır.
Parçacık ve spin dinamiği, 4.~derece Gauss--Legendre (GL4) simplektik
entegratörle, Thomas--BMT denkleminin tam (yaklaşıksız) formu
kullanılarak izlenir; EDM terimi $\eta$ anahtarıyla açılıp
kapatılabilir. \todo{adım boyu, tur sayısı, $t_2$ değerleri tablosu}

Kaçıklıklar, her kuadrupolün düşey konumuna eklenen 48 bileşenli
$\boldsymbol{\delta y}$ vektörüyle temsil edilir. Analiz, bu vektörün
FODO-periyodik Fourier tabanındaki bileşenleri üzerinden yürür:
$k$ modu, hücre indeksi $i$ üzerinde $\cos(2\pi k i/24)$ ve
$\sin(2\pi k i/24)$ ağırlıklarıyla tanımlanır; QF ve QD'nin zıt
işaretli (antisimetrik) ve aynı işaretli (simetrik) kombinasyonları
ayrı alt-uzaylar oluşturur (Bölüm~\ref{sec:taban}).

\subsection{Sahte EDM mekanizması ve doğrusal model}

Kaçıklık deseni $\boldsymbol{\delta y}$, kapalı yörünge bozulması
üzerinden net bir radyal alan örneklemesine ve dolayısıyla
\begin{equation}
f \equiv \frac{dS_y}{dt}
= \sum_k \left( c_k^{\cos} a_k^{\cos} + c_k^{\sin} a_k^{\sin} \right)
\label{eq:dogrusal-model}
\end{equation}
biçiminde, mod genlikleri $a_k$'da doğrusal bir sahte EDM hızına yol
açar. $c_k$ katsayıları kafesin geometrik özelliğidir: radyal alanın
yörünge boyunca ağırlıklı integrali olduğundan \emph{rezonant
değildir} ve hiçbir $k$ için sıfırlanmaz (ölçülen değerler:
$|c_2|=90.7$, $|c_3|=23.6$, $|c_4|=8.4$, $|c_7|\approx2.0$~rad/s/m
\todo{tam tablo}). Bu, betatron rezonansıyla şekillenen yörünge
kazançlarından (Bölüm~\ref{sec:orbit-trim}) temelden farklı bir
davranıştır ve iki kademenin farklı modları görmesinin kökenidir.

[SEKIL 1: $c_k$ haritası $k=1..24$, üç arka planda evrensellik]

%==============================================================
\section{Spin kademesi: ölç--trimle döngüsü}
\label{sec:spin-trim}

\subsection{Üç önkoşulun kanıtı}

Spin-sürülü trimin işlerliği üç özelliğe dayanır; üçü de bilinen
desenlerle kontrollü simülasyonlarda sınanmıştır:

\textbf{(i) Doğrusallık.} Tek modun genliği taranırken $f$ genlikle
orantılı kalır; $\SI{100}{\micro\meter}$ mertebesine kadar ölçülen
sapma $\%0.0$'dır. \todo{tarama şekli}

\textbf{(ii) Desenden bağımsızlık.} $c_k$, üzerine bindiği arka plan
deseninden bağımsızdır: sıfır, orta ve güçlü rastgele arka planlarda
ölçülen $c_k$ haritaları arasındaki korelasyon $1.0000$'dır. Bu,
kalibrasyonun (bilinmeyen) gerçek desen mevcutken yapılabilmesini
sağlar --- yöntemin saha uygulanabilirliğinin temeli.

\textbf{(iii) Fırlatma koşulundan bağımsızlık.} Betatron salınımı
yapan tekil parçacıklarda $c_k$, demet ortalaması alındığında kapalı
yörünge değerine $\%0.5$ içinde döner; trim, demet polarimetre
sinyaliyle çalışabilir.

\subsection{Faz problemi ve çift-kuadratür kalibrasyon}

Gerçek desenin her modu bilinmeyen bir fazla gelir:
$a_k^{\cos},a_k^{\sin}$ çifti. Gözlenebilir tek skaler olduğundan
($f$), her modun iki kuadratüründe birer kalibrasyon vuruşu
($+A_\mathrm{cal}$ cos, $+A_\mathrm{cal}$ sin) yapılır; iki ölçüm
modun hem genliğini hem fazını çözer. $k=1..N_k$ için toplam
$2N_k+1$ dolum yeterlidir.

\subsection{Performans}

Bilinen rastgele desenlerde ($\sigma_q=\SI{100}{\micro\meter}$ RMS)
tam kalibrasyonlu tek atış trimi $f$'yi $2\times10^{7}$ kat bastırır;
kalibrasyonun artıklarıyla yürüyen iteratif ölç--trimle döngüsü desen
bilgisi olmadan $\sim10^{3}\times$ bastırmaya birkaç turda ulaşır.
BPM gürültüsü ve ofseti bu kademeyi etkilemez (BPM kullanılmaz);
pratik sınır polarimetre istatistiğidir. \todo{polarimetre istatistik
modeli ve süre bütçesi}

[SEKIL 2: iteratif döngü yakınsaması; tek atış histogramı]

\subsection{EDM güvenliği: spin kademesi neden tek başına kullanılamaz}
\label{sec:edm-guvenlik}

$f$ gözlenebiliri gerçek EDM katkısını da içerir. Döngü körlemesine
$f\to0$ hedeflerse gerçek sinyali de trimin içine gömer. Güvenli
kullanım iki yoldan biriyle olur: (a)~CW ve CCW saklamaların farkı
üzerinden --- sahte EDM yön değiştirince işaret değiştirir, gerçek EDM
değiştirmez --- yalnız-sistematik bir gözlenebilir kurmak;
(b)~kademeyi yalnızca gerçek EDM duyarlılığının çok altındaki kaba
artıkları temizlemek için, sınırlı çevrimle çalıştırmak. Bu
gereklilik, spin kademesinin önüne hızlı ve yapısal olarak EDM-kör
bir kaba kademe koymanın ana gerekçesidir.

%==============================================================
\section{Yörünge kademesi: EDM-kör kaba trim}
\label{sec:orbit-trim}

\subsection{Yöntem}

Her mod düğmesinin yörünge parmak izi, $A_\mathrm{cal}$ genlikli
kontrollü vuruşlarla diferansiyel olarak kalibre edilir
($\mathbf{o}_k=(\mathbf{y}(A_\mathrm{cal})-\mathbf{y}_0)/A_\mathrm{cal}$);
statik BPM ofseti farkta düşer, kalibrasyon matrisine girmez. Ölçülen
yörünge $\mathbf{y}$, $O=[\mathbf{o}_1\ldots]$ matrisine en küçük
karelerle fit edilir; kestirilen genlikler ters işaretle mekanik
kaydırıcılara uygulanır. Spin hiçbir aşamada okunmaz.

\subsection{Kazanç yasası ve fit eşiği}

Modların yörünge kazançları $G_k=\|\mathbf{o}_k\|_\mathrm{rms}$
betatron rezonansıyla şekillenir ve klasik düz-yaklaşım kapalı-yörünge
yanıtına oturur:
\begin{equation}
G_k=\frac{C}{|Q_\mathrm{eff}^2-k^2|},\qquad C=24.8,\;
Q_\mathrm{eff}^2=5.03
\label{eq:kazanc}
\end{equation}
($k\le6$ ölçümlerine RMS $\%0.6$ ile oturtulmuştur). Yasanın
\emph{öngörü} gücü, fit aralığının tamamen dışındaki $k=7..12$
modlarının sonradan kalibre edilmesiyle sınanmıştır: $k\le9$'da sapma
$\le\%1$. [SEKIL 3: kazanç yasası, ölçüm + öngörü]

Statik BPM ofseti kalibrasyona girmese de tek-atış fit'e girer ve mod
başına $\varepsilon_k=\sigma_b\sqrt{2/N}/G_k$ yanlılığı bırakır.
Modun beklenen gerçek içeriği $A_k\approx\sigma_q\sqrt{2/N}$
olduğundan, modu fit etmenin kazandırma koşulu
\begin{equation}
G_k>\frac{\sigma_b}{\sigma_q}
\;\;\Longleftrightarrow\;\;
k_\mathrm{max}^2<Q_\mathrm{eff}^2+C\,\frac{\sigma_q}{\sigma_b}
\label{eq:esik}
\end{equation}
kapalı biçimini alır. $\sigma_b=\sigma_q=\SI{100}{\micro\meter}$ için
$k_\mathrm{max}=5.5$: $k=4$ güvenli, $k=5$ tam eşikte (seed'e bağlı
kararsız davranışın nedeni), $k=6$ zararlı --- dört fit-genişliği
varyantının simülasyon sonuçlarıyla birebir uyumlu.
\todo{varyant tablosu: A/C/D/B, seed=321 ve 5-seed RMS}

Gerçek makinede bağlayıcı olan simülasyon sayıları değil bu
prosedürdür: $G_k$ devreye almada aynı diferansiyel kalibrasyonla
yerinde ölçülür, eşik o makinenin $\sigma_b$ değeriyle yeniden
hesaplanır.

\subsection{Mod ayrışımı: korelasyon değil, eşik}

Fit kesiminin modlar-arası korelasyondan kaynaklanmadığı doğrudan
ölçülmüştür: 24 düğmenin normalize yörünge parmak izlerinin Gram
matrisinde $k\le6$ bloğunun en büyük köşegen-dışı elemanı $0.011$;
fit-dışı antisimetrik modların kestirimlere sızıntısı
$\max|L|=0.002$. Yörünge fit'i modları tek atışta, neredeyse dik
biçimde ayırır --- spin gözlenebilirinin rank-1 yapısıyla (tüm modlar
tek skalere çöker; tekil mod bastırma $|f|$'yi artırabilir) keskin
karşıtlık içindedir. [SEKIL 4: Gram matrisi]

%==============================================================
\section{Taban analizi: yörüngenin göremediği yarı ve spinin gerekliliği}
\label{sec:taban}

Yörünge kademesinin seed'ler arası artığı $\sim2.5\times10^{-4}$~rad/s
RMS'de doyar --- zincir hedefi $10^{-5}$'in $\sim25$ katı. Kaynağı,
kaçıklık uzayının dik ayrıştırmasıyla kanıtlanmıştır.

48 serbestlik derecesinin antisimetrik Fourier bazında ($k=0..12$,
QF/QD zıt işaretli) yaşayan kısmı 25 boyuttur; kalan 23 boyut
\emph{simetrik} (QF ve QD aynı yönde) kombinasyonlardır. Zıt gradyan
işaretleri nedeniyle simetrik bir kaçıklıkta QF ve QD'nin yörünge ve
spin tekmeleri birbirini söndürmeye çalışır; söndürme iki kanalda da
\emph{kısmidir}. Seed=321 deseninin iki dik parçası ayrı ayrı
simüle edildiğinde:

\begin{table}[h]
\centering
\caption{Dik ayrıştırma: antisimetrik ve simetrik parçaların ayrı ayrı
ürettiği yörünge ve sahte EDM ($\sigma_q=\SI{100}{\micro\meter}$,
seed=321; parçaların iç çarpımı $\sim10^{-24}$).}
\begin{tabular}{lccc}
\toprule
parça & desen RMS [\si{\micro\meter}] & COD RMS [\si{\micro\meter}]
& $dS_y/dt$ [rad/s] \\
\midrule
tam            & 89.5 & 628 & $-1.62\times10^{-3}$ \\
antisim (25 b.) & 69.8 & 550 & $-1.52\times10^{-3}$ \\
simetrik (23 b.) & 54.9 & 165 & $-1.02\times10^{-4}$ \\
\bottomrule
\end{tabular}
\label{tab:ayristirma}
\end{table}

Birim kaçıklık başına: yörünge kazancı antisimetrikte 7.9, simetrikte
3.0 ($\sim2.6\times$ iptal); spin kuplajı antisimetrikte
$2.2\times10^{-5}$, simetrikte $1.9\times10^{-6}$~rad/s/\si{\micro\meter}
($\sim12\times$ iptal). Spindeki kalıntının kaynağı, QF--QD arasındaki
deflektör arklarında biriken geometrik spin fazıdır
($a\gamma\,\Phi_\mathrm{def}\approx0.23$~rad/yarı-hücre). Simetrik
parçanın tek başına taşıdığı $1.0\times10^{-4}$~rad/s, ölçülen yörünge
trimi tabanının ana bileşenidir.

\textbf{Simetrik modlar fit'e eklenemez.} Simetrik içerik yörüngede
görünür olsa da kazançları ($0.7$--$4.7$) eşik
(Denk.~\ref{eq:esik}) civarında veya altındadır. Genişletilmiş bazla
($16$ sütun) yapılan fit denemesi bunu doğrulamıştır: BPM ofseti
kestirimlere enjekte olur (kestirilen simetrik genlikler
$\SI{136}{\micro\meter}$ RMS; gerçek içerik $\SI{55}{\micro\meter}$)
ve artık $9\times$ \emph{kötüleşir}. Fit edilebilirliği belirleyen
korelasyon değil kazançtır; eşik altındaki bilgi BPM'den
çıkarılamaz.

Sonuç: yörünge kademesi $\sigma_b=\SI{100}{\micro\meter}$'de
$\sim10^{-4}$~rad/s tabanına demirler. Bu tabanı yalnız, tüm 48
boyutu gören spin kademesi (hiçbir modda $c_k=0$ değil ---
simetrik modların kuplajı bastırılmış ama sonlu) aşabilir. İki kademe
böylece zorunlu olarak tamamlayıcıdır.

[SEKIL 5: ayrıştırma sonuçları; genişletilmiş baz karşılaştırması]

%==============================================================
\section{Elenen alternatifler}
\label{sec:elenen}

Aşağıdaki seçenekler denenmiş ve mekanizmaları anlaşılarak elenmiştir;
yöntemin sınırlarını çizdikleri için raporlanmaktadır.

\textbf{Tekil mod spin trimi.} $f_0=\sum c_k a_k$ toplamında terimler
kısmen iptal hâlinde olabilir; yalnız $k=2$'yi bastırmak kalan toplamı
büyütebilir. Spin gözlenebiliri rank-1 olduğundan modlar ancak
çok-modlu, fazlı fit ile birlikte trimlenebilir.

\textbf{COD minimizasyonu.} $N=2..24$ korrektörlü klasik yörünge
düzeltmesi COD'yi $\%50$--$95$ azaltırken sahte EDM yalnız $\sim\%15$
düşer: yörünge ve $f$, kaçıklık uzayının farklı izdüşümleridir.
Yörünge düzeltmesi gerekli ama yeterli değildir.

\textbf{Genişletilmiş (simetrik) fit bazı.} Bölüm~\ref{sec:taban}:
eşik altı sütunlar ofset enjeksiyonuyla sonucu $9\times$ kötüleştirir.

\textbf{Radyal başlangıç polarizasyonu.} Radyal spin
($S_x=1$, $S_z=0$) hem gerçek hem sahte EDM dönüşüne birinci mertebede
kördür (ölçüm: EDM sinyali $3.5\times10^{6}\times$ bastırılır);
kaçıklık sinyalinin $\%5$'i $\Omega_z$ kanalında kalır. EDM'siz bir
sistematik kanalı olarak değerlidir; fakat $\Omega_z$ düşey yörünge
eğimiyle sürüldüğünden, yörünge üretmeyen simetrik içeriğe o da
büyük ölçüde kördür --- taban problemini çözmez.

\textbf{Kafes topolojisi değişiklikleri.} Deflektörü kuadrupol
çiftinin dışına alan hücre (QF--drift--QD--drift--ARC) simetrik spin
iptalini onarabilirdi; ince-mercek analizi aynı ton için $\sim2\times$
kuadrupol gücü gerektiğini ve tam bir kafes yeniden tasarımı
anlamına geldiğini göstermiştir. Tek tip kuadrupol (yalnız QD)
seçeneği ise quad-flip iptal simetrisini --- $10^{-5}\to10^{-9}$
kademesinin temelini --- yok ettiğinden elenmiştir.

%==============================================================
\section{Deneysel reçete ve süre bütçesi}
\label{sec:recete}

\begin{enumerate}
\item \textbf{Yörünge kalibrasyonu:} $k=1..4$ cos/sin düğmeleri
diferansiyel vuruşlarla kalibre edilir (8 yörünge ölçümü, dakikalar).
$G_k$ yerinde ölçülür; Denk.~(\ref{eq:esik}) ile fit aralığı o
makinenin $\sigma_b$'siyle belirlenir.
\item \textbf{Yörünge trimi:} Tur-ortalamalı BPM okuması, LSQ fit,
mekanik geri besleme. Beklenen sonuç: antisimetrik $k\le
k_\mathrm{max}$ içeriği birkaç \si{\micro\meter}'ye iner; artık
$\sim10^{-4}$~rad/s.
\item \textbf{Spin kalibrasyonu ve trimi:} CW/CCW farkı üzerinden
yalnız-sistematik gözlenebilirle $c_k$ çift-kuadratür kalibrasyonu
($2N_k+1$ dolum) ve çok-modlu trim; hedef $10^{-5}$~rad/s altı.
\item \textbf{İptal zinciri:} Kalan artık CW/CCW + quad-flip
teknikleriyle $10^{-9}$~rad/s altına taşınır~\cite{todo-omarov}.
\end{enumerate}

\todo{dolum sayısı/süre bütçesi tablosu; polarimetre istatistiği ile
spin kademesi çözünürlüğünün bağlanması}

%==============================================================
\section{Tartışma ve sonraki çalışmalar}
\label{sec:tartisma}

\todo{Bu bölüm genişletilecek. İskelet:}

(i)~Tüm sonuçlar tek kafes ve tek enerji içindir; eşik prosedürünün
makine-bağımsızlığı iddiası ikinci bir çalışma noktasıyla (farklı ton
veya hücre sayısı) desteklenecektir. (ii)~RF ve sekstüpoller açıkken
doğrusallık teyidi. (iii)~$t_2=\SI{1}{ms}$ izleme süresinden spin
koherans sürelerine ekstrapolasyon. (iv)~Yatay kaçıklıklar ve kuadrupol
tilt çapraz terimleri. (v)~Zamanla sürüklenen hizalamada sürekli
izleme kipi: yörünge kademesi saniyeler ölçeğinde tekrarlanabilir
olduğundan run sırasında drift takibine uygundur.

%==============================================================
\section{Sonuç}
\label{sec:sonuc}

Sahte EDM, hizalama vektörü hiç ölçülmeden, iki tamamlayıcı
geri-besleme kademesiyle kaynağında bastırılabilir. EDM-kör yörünge
kademesi, kapalı-form eşik teorisiyle ($G_k$ yasası +
$k_\mathrm{max}$ koşulu) hangi modları güvenle fit edebileceğini
önceden bilir ve antisimetrik içeriği dakikalar içinde temizler;
tabanını oluşturan simetrik içerik ise yalnız spin kademesinin
görebildiği kanaldadır ve çift-kuadratür kalibrasyonlu ölç--trimle
döngüsüyle temizlenir. Mimari, CW/CCW + quad-flip iptal zincirinin
$10^{-5}$~rad/s giriş gereksinimine giden bütünsel ve deneysel olarak
uygulanabilir bir yol sunmaktadır.

\section*{Teşekkür}
\todo{teşekkür ve fon bilgisi}

\begin{thebibliography}{9}
\bibitem{todo-pedm} \todo{pEDM tasarım raporu / öneri makalesi}
\bibitem{todo-omarov} \todo{Omarov ve ark., PRD --- tolerans ve
CW/CCW + quad-flip iptal zinciri}
\bibitem{todo-bba} \todo{BBA / k-modülasyon referansları}
\bibitem{todo-loco} \todo{LOCO / tepki matrisi referansı}
\end{thebibliography}

\end{document}
